% SEGMENT OLP-0637-S01
% Part: history
% Chapter: set-theory
% Section: mythology

\documentclass[../../../include/open-logic-section]{subfiles}

\begin{document}

\olfileid{his}{set}{mythology}

\olsection{வரலாற்றைவிடத் தொன்மக் கதை அதிகமா?}

ஒரு நூற்றாண்டுக்கும் மேலான இடைவெளிக்குப் பிறகு இந்த நிகழ்வுகளைத்
திரும்பிப் பார்க்கும்போது, முன்பு கூறப்பட்ட முடிவுகளை அப்படியே
நம்பிவிடக் கூடாது. அந்தக் கண்டுபிடிப்புகள் புதுமையானவை,
ஆர்வமூட்டுபவை, வியப்பளிப்பவை என்பதில் ஐயமில்லை. ஆனால் அவை
உண்மையில் எந்த அளவுக்கு அதிர்ச்சியளித்தன? வடிவியல் உள்ளுணர்வை
நம்பக் கூடாது என்பதை அவை உண்மையிலேயே காட்டினவா?

அதிர்ச்சி பற்றிப் பேசும்போது, காண்டோர் டெடெகிண்டிற்கு எழுதிய
“\emph{je le vois, mais je ne le crois
pas}” என்ற புகழ்பெற்ற வாசகம் அதன் சூழலிலிருந்து பெருமளவு
பிரித்தெடுக்கப்படுகிறது என்று \citet{Gouvea2011} சுட்டிக்காட்டுகிறார்.
அந்தச் சூழலின் கூடுதல் பகுதி இதோ (\citeauthor{Gouvea2011}-இலிருந்து
மேற்கோள்):
\begin{quote}
இந்தப் பொருளின்மீது எனக்குள்ள தீவிர ஈடுபாட்டால் உங்களுடைய
அன்பையும் பொறுமையையும் இவ்வளவு கோர வேண்டியிருப்பதை மன்னியுங்கள்.
அண்மையில் நான் உங்களுக்கு அனுப்பிய செய்திகளே எனக்கும்கூட
எதிர்பாராதவையாகவும் புதுமையானவையாகவும் உள்ளன. மதிப்பிற்குரிய
நண்பரே, அவை சரியா என்பது பற்றிய உங்கள் முடிவைப் பெறும்வரை
என் மனம் அமைதியடையாது. நீங்கள் என்னோடு உடன்படாதவரையில்
என்னால் சொல்ல முடிவது இதுதான்:
\emph{je le vois, mais je ne le crois pas.}
\end{quote}
தனது முடிவு “மிகவும் எதிர்பாராதது; மிகவும் புதுமையானது” என்பதை
காண்டோர் அறிந்திருந்தார். ஆனால் அதை அவர் எப்போதுமே
\emph{நம்பமுடியாதது} என்று கருதினார் என்பது சந்தேகமே.
\citeauthor{Gouvea2011} சுட்டிக்காட்டுவதுபோல், தாம் அளித்த
நிறுவலை டெடெகிண்ட் சரிபார்க்க வேண்டும் என்றுதான் அவர் கேட்டார்.

% SEGMENT OLP-0637-S02
இப்போது வடிவியல் உள்ளுணர்வைப் பார்ப்போம். பியானோ தனது
வெளிநிரப்பு வளைவை எந்த வரைபடமும் சேர்க்காமல் வெளியிட்டார்.
ஆனால் ஹில்பர்ட் தனது வளைவை வெளியிட்டபோது, தன் நோக்கத்தை
விளக்கினார்: வாசகர்கள் “பின்வரும் வடிவியல் உள்ளுணர்வின்
உதவியை நாடினால்”, பியானோவின் முடிவைத் தெளிவாகப் புரிந்துகொள்ள
வழி கிடைக்கும் என்றார். பின்னர்
\olref[his][set][pathology]{sec} பகுதியில் இருப்பவற்றைப்
போன்ற \emph{வரைபடங்கள்} பலவற்றைச் சேர்த்தார்.

பொதுவாக, வெளியிடப்படும் நிறுவல்களில் வரைபடங்கள் ஓரளவு
புழக்கத்திலிருந்து விலகியிருந்தாலும், கணிதவியலாளர்கள்
நிறுவும்போது அவற்றை \emph{அடிக்கடி} பயன்படுத்துகிறார்கள்
என்ற உண்மையை மறுக்க முடியாது. (சுருக்கமாகச் சொன்னால்,
வடிவியல் உள்ளுணர்வை எப்போது நம்பலாம் என்பதை நல்ல
கணிதவியலாளர்கள் அறிவார்கள்.)

எனவே பரபரப்பாகச் சொல்லப்படும் கதையை உடனே நம்ப வேண்டாம்;
குறைந்தது, அதைச் சரிபார்க்காமல் ஏற்க வேண்டாம்.
மேலும் அறிய \citet{Giaquinto2007}-ஐப் படிக்கலாம்.

\end{document}
